\documentclass[11pt]{article}

\usepackage[T1]{fontenc}
\usepackage[utf8]{inputenc}
\usepackage{lmodern}
\usepackage[margin=1in]{geometry}
\usepackage{microtype}
\usepackage{amsmath,amssymb,amsthm,mathtools,mathrsfs}
\usepackage{enumitem}
\usepackage[colorlinks=true,linkcolor=blue,urlcolor=blue]{hyperref}
\setlist[enumerate]{leftmargin=*,itemsep=4pt,topsep=4pt}

\newtheorem{theorem}{Theorem}[section]
\newtheorem{lemma}[theorem]{Lemma}
\newtheorem{proposition}[theorem]{Proposition}
\newtheorem{corollary}[theorem]{Corollary}
\theoremstyle{remark}
\newtheorem{remark}[theorem]{Remark}

\newcommand{\R}{\mathbb R}
\newcommand{\C}{\mathbb C}
\newcommand{\Z}{\mathbb Z}
\newcommand{\T}{\mathbb T}
\newcommand{\F}{\mathcal F}
\newcommand{\Sstar}{\mathcal S_\star}
\newcommand{\abs}[1]{\left\lvert#1\right\rvert}
\newcommand{\norm}[1]{\left\lVert#1\right\rVert}
\newcommand{\dd}{\,\mathrm d}
\DeclareMathOperator*{\argmax}{arg\,max}
\DeclareMathOperator{\supp}{supp}

\title{E16: Linear Spectral Concentration and\\
Convergence of the Dominant Fourier Wavenumber}
\author{Manuscript 1 --- Self-Emergent Fourier Cymatics\\
\small A self-contained proof with explicit hypotheses and counterexample}
\date{October 26, 2026}

\begin{document}
\maketitle

\begin{abstract}
For the linearized finite-band evolution
$\partial_t A=(r-a\Delta-b\Delta^2)A$ on $\R^d$, with $a,b>0$,
we establish the exact spectral-power law
$P(\xi,t)=P(\xi,0)e^{2\lambda(\xi)t}$, where
$\lambda(\xi)=r+a\abs{\xi}^2-b\abs{\xi}^4$.
Its unique maximizing \emph{radius} is
$k_\star=\sqrt{a/(2b)}$. If the initial Fourier power has positive
mass in every neighborhood of $\{\abs\xi=k_\star\}$, the normalized
spectral measure concentrates exponentially in any fixed tubular
neighborhood of that shell. Under an additional continuous,
vanishing-at-infinity hypothesis on the initial spectral power,
every maximizing Fourier wavenumber tends radially to $k_\star$.
We also give a quantitative $O(t^{-1/2})$ bound when a shell frequency
is initially populated, a radial-spectrum formulation, a torus
qualification, and a counterexample showing why the initial-support
hypothesis is necessary. The result is linear spectral selection,
not a nonlinear localization or confinement theorem.
\end{abstract}

\section{Model, notation, and hypotheses}
Fix a spatial dimension $d\ge 1$, $a,b>0$, and $r\in\R$.
Consider
\begin{equation}\label{eq:flow}
 \partial_t A=(r-a\Delta-b\Delta^2)A,
 \qquad A(0)=A_0\in L^2(\R^d;\C),\quad A_0\ne0.
\end{equation}
The unitary Fourier transform is
\begin{equation}\label{eq:fourier}
 \widehat A(\xi)=(2\pi)^{-d/2}
 \int_{\R^d}e^{-ix\cdot\xi}A(x)\dd x,
\end{equation}
initially on Schwartz functions and then extended by Plancherel.
Put
\begin{align}
 k_\star&=\sqrt{\frac{a}{2b}}>0,
 &\mu&=r+\frac{a^2}{4b}, \label{eq:params}\\
 \lambda(\xi)&=r+a\abs{\xi}^2-b\abs{\xi}^4
 =\mu-b\bigl(\abs{\xi}^2-k_\star^2\bigr)^2,
 &q(\xi)&=\bigl(\abs{\xi}^2-k_\star^2\bigr)^2.
 \label{eq:lambda}
\end{align}
The preferred shell is
\[
 \Sstar=\{\xi\in\R^d:\abs{\xi}=k_\star\}.
\]
The relevant spectral-support hypothesis is
\begin{equation}\label{eq:H}
 \tag{H}
 \int_{\{\xi:\,\lvert\abs\xi-k_\star\rvert<\varepsilon\}}
 \abs{\widehat A_0(\xi)}^2\dd\xi>0
 \qquad\text{for every }\varepsilon>0.
\end{equation}
This is an \emph{essential-support} condition: it does not require
nonzero pointwise power on the measure-zero sphere itself.

\section{Exact Fourier evolution}
\begin{proposition}[Exact modal power law]\label{prop:exact}
For every $A_0\in L^2(\R^d)$, equation \eqref{eq:flow} has a
unique mild solution $A\in C([0,\infty);L^2(\R^d))$, and
\begin{align}
 \widehat A(t,\xi)&=e^{\lambda(\xi)t}\widehat A_0(\xi),
 \label{eq:ahat}\\
 P(\xi,t):=\abs{\widehat A(t,\xi)}^2
 &=P(\xi,0)e^{2\lambda(\xi)t}
 \quad\text{for a.e. }\xi\in\R^d.
 \label{eq:power}
\end{align}
\end{proposition}
\begin{proof}
Using
$\widehat{\Delta A}=-\abs\xi^2\widehat A$ and
$\widehat{\Delta^2 A}=\abs\xi^4\widehat A$,
Fourier transformation gives, frequency by frequency,
$\partial_t\widehat A=\lambda(\xi)\widehat A$.
Solving this scalar ODE yields \eqref{eq:ahat}.
Since $\lambda(\xi)\le\mu$,
$\abs{e^{t\lambda(\xi)}}\le e^{t\mu}$ for $t\ge0$,
so multiplication by $e^{t\lambda}$ is bounded on $L^2$.
The semigroup law is immediate, and dominated convergence yields
strong continuity at $t=0$ (and hence at all $t\ge0$).
The Fourier multiplier therefore defines the unique $L^2$ mild
solution generated by $r-a\Delta-b\Delta^2$.
Taking absolute values squared proves \eqref{eq:power}.
\end{proof}

\begin{lemma}[Unique maximizing radius and strict gap]\label{lem:gap}
The set of Fourier vectors maximizing $\lambda$ is exactly $\Sstar$.
For every $\delta>0$,
\begin{equation}\label{eq:gap}
 \bigl\lvert\abs\xi-k_\star\bigr\rvert\ge\delta
 \quad\Longrightarrow\quad
 \mu-\lambda(\xi)\ge bk_\star^2\delta^2.
\end{equation}
\end{lemma}
\begin{proof}
Completion of the square in \eqref{eq:lambda} shows
$\lambda\le\mu$, with equality exactly when
$\abs\xi=k_\star$. Furthermore,
\[
 \mu-\lambda(\xi)
 =b(\abs\xi-k_\star)^2(\abs\xi+k_\star)^2
 \ge bk_\star^2(\abs\xi-k_\star)^2.
\]
This proves \eqref{eq:gap}.
\end{proof}

\section{Concentration of normalized spectral power}
The integral $Z(t):=\int_{\R^d}P(\xi,t)\dd\xi$ is finite and
strictly positive for every finite $t\ge0$. Define the probability
measure
\begin{equation}\label{eq:nu}
 \nu_t(E)=\frac{\int_E P(\xi,t)\dd\xi}{Z(t)}
 \quad\text{for measurable }E\subset\R^d.
\end{equation}

\begin{theorem}[E16: concentration on the preferred shell]
\label{thm:concentration}
Suppose $A_0\ne0$ belongs to $L^2(\R^d)$ and satisfies
\eqref{eq:H}. For every $\delta>0$, let
\[
 E_\delta=\{\xi:\lvert\abs\xi-k_\star\rvert\ge\delta\}.
\]
Then there exists $C_\delta\in(0,\infty)$, depending only on the
initial spectrum and $\delta$, such that
\begin{equation}\label{eq:concentration}
 \boxed{\displaystyle
 \nu_t(E_\delta)
 \le C_\delta e^{-bk_\star^2\delta^2t}
 \xrightarrow[t\to\infty]{}0.}
\end{equation}
Equivalently, for every $\delta>0$,
\[
 \lim_{t\to\infty}
 \nu_t\bigl(\{\xi:\lvert\abs\xi-k_\star\rvert<\delta\}\bigr)
 =1.
\]
\end{theorem}
\begin{proof}
Write
$f(\xi)=\abs{\widehat A_0(\xi)}^2\in L^1(\R^d)$,
$F=\int_{\R^d}f\dd\xi=\norm{A_0}_{L^2}^2>0$, and
$m_\delta=k_\star^2\delta^2>0$.
By \eqref{eq:lambda} and \eqref{eq:power},
\begin{equation}\label{eq:P-factor}
 P(\xi,t)=e^{2\mu t}f(\xi)e^{-2btq(\xi)}.
\end{equation}
On $E_\delta$ we have $q\ge m_\delta$ by
Lemma~\ref{lem:gap}, so
\begin{equation}\label{eq:outer}
 \int_{E_\delta}P(\xi,t)\dd\xi
 \le e^{2\mu t}e^{-2bm_\delta t}F.
\end{equation}
Consider the open inner region
\[
 B_\delta=\{\xi:q(\xi)<m_\delta/2\}.
\]
Because $q$ is continuous and vanishes on $\Sstar$,
$B_\delta$ contains some tubular neighborhood
$\{\lvert\abs\xi-k_\star\rvert<\varepsilon_\delta\}$.
Hypothesis \eqref{eq:H} therefore implies
\[
 c_\delta:=\int_{B_\delta}f(\xi)\dd\xi>0.
\]
For $\xi\in B_\delta$, $e^{-2btq(\xi)}\ge e^{-bm_\delta t}$,
whence
\begin{equation}\label{eq:inner}
 Z(t)\ge e^{2\mu t}e^{-bm_\delta t}c_\delta.
\end{equation}
Dividing \eqref{eq:outer} by \eqref{eq:inner},
\[
 \nu_t(E_\delta)
 \le\frac{F}{c_\delta}e^{-bm_\delta t}.
\]
Taking $C_\delta=F/c_\delta$ and
$m_\delta=k_\star^2\delta^2$ proves the result.
\end{proof}

\begin{remark}[Normalized does not mean persistent amplitude]
The theorem concerns $P/Z$, not the magnitude of $Z(t)$.
In particular, if $\mu=0$, then $P(\xi,t)\to0$ for almost
every $\xi$, since the exact shell has Lebesgue measure zero;
dominated convergence gives $\norm{A(t)}_2\to0$.
Nevertheless $\nu_t$ still concentrates on the shell under (H).
\end{remark}

\section{Convergence of the spectral argmax}
For $L^2$ data alone, the Fourier transform and its pointwise
maximum are not intrinsically defined. The next theorem states the
additional regularity needed for a pointwise argmax.

\begin{theorem}[Argmax convergence]\label{thm:argmax}
In addition to the hypotheses of Theorem~\ref{thm:concentration},
assume that the initial power
$f(\xi)=\abs{\widehat A_0(\xi)}^2$ has a nonnegative continuous
representative in $C_0(\R^d)$ (continuous and vanishing at infinity).
For every $t\ge0$, the set $\argmax_{\xi\in\R^d}P(\xi,t)$ is
nonempty. Every choice
$\xi_t\in\argmax P(\cdot,t)$ satisfies
\begin{equation}\label{eq:argmax}
 \boxed{\displaystyle
 \lim_{t\to\infty}\abs{\xi_t}=k_\star.}
\end{equation}
This is radial convergence and does not select one vector direction.
\end{theorem}
\begin{proof}
Since $f\in C_0$, $f\ge0$, and $f\not\equiv0$,
$M:=\sup_{\xi}f(\xi)$ lies in $(0,\infty)$.
For each finite $t$, \eqref{eq:P-factor} defines a nonzero
nonnegative function in $C_0(\R^d)$; hence its positive global
maximum is attained.

Fix $\delta>0$ and set
$m_\delta=k_\star^2\delta^2$.
Since (H) holds and $q$ vanishes continuously on the preferred
shell, there exists $\eta_\delta\in\R^d$ such that
\begin{equation}\label{eq:eta}
 f(\eta_\delta)>0,
 \qquad q(\eta_\delta)<m_\delta/2.
\end{equation}
At this point,
\begin{equation}\label{eq:comparison-in}
 P(\eta_\delta,t)
 \ge e^{2\mu t}f(\eta_\delta)e^{-bm_\delta t}.
\end{equation}
For every $\xi\in E_\delta$,
\begin{equation}\label{eq:comparison-out}
 P(\xi,t)\le e^{2\mu t}M e^{-2bm_\delta t}.
\end{equation}
Combining the bounds yields
\begin{equation}\label{eq:ratio}
 \frac{\sup_{\xi\in E_\delta}P(\xi,t)}
 {P(\eta_\delta,t)}
 \le\frac{M}{f(\eta_\delta)}e^{-bm_\delta t}
 \xrightarrow[t\to\infty]{}0.
\end{equation}
Thus for all sufficiently large $t$ the power at $\eta_\delta$
exceeds the power at every $\xi\in E_\delta$.
Every global maximizer $\xi_t$ then lies outside $E_\delta$, i.e.,
$\lvert\abs{\xi_t}-k_\star\rvert<\delta$.
As $\delta>0$ was arbitrary, \eqref{eq:argmax} follows.
\end{proof}

\begin{corollary}[Explicit $t^{-1/2}$ upper bound]\label{cor:rate}
Under the hypotheses of Theorem~\ref{thm:argmax}, suppose there is
$\eta_\star\in\Sstar$ with $f(\eta_\star)>0$.
Let $M=\sup_\xi f(\xi)$. Then, for every $t>0$ and every
$\xi_t\in\argmax P(\cdot,t)$,
\begin{equation}\label{eq:rate}
 \boxed{\displaystyle
 \bigl\lvert\abs{\xi_t}-k_\star\bigr\rvert
 \le\sqrt{\frac{\log(M/f(\eta_\star))}
 {2bk_\star^2t}}.}
\end{equation}
In particular the radial argmax error is $O(t^{-1/2})$.
\end{corollary}
\begin{proof}
Because $q(\eta_\star)=0$ and $\xi_t$ maximizes $P$,
\[
 f(\xi_t)e^{-2btq(\xi_t)}
 \ge f(\eta_\star).
\]
Since $f(\xi_t)\le M$,
\[
 q(\xi_t)\le\frac{\log(M/f(\eta_\star))}{2bt}.
\]
The factorization
$q(\xi_t)=(\abs{\xi_t}-k_\star)^2(\abs{\xi_t}+k_\star)^2$
and $\abs{\xi_t}+k_\star\ge k_\star$ give \eqref{eq:rate}.
\end{proof}

\section{Radial spectral density: the scalar wavenumber formulation}
In dimension $d>1$, a Fourier frequency is a vector $\xi$;
$k$ usually denotes its nonnegative radius. Define, for a.e. $k>0$,
\begin{equation}\label{eq:radial}
 \mathscr P(k,t)
 =k^{d-1}\int_{\mathbb S^{d-1}}P(k\omega,t)
 \dd\omega.
\end{equation}
For $d=1$, the integral is over the two-point sphere
$\mathbb S^0=\{-1,+1\}$ with counting measure.
Polar integration gives
$\int_0^\infty\mathscr P(k,t)\dd k=Z(t)$.
Because $\lambda(k\omega)=\mu-b(k^2-k_\star^2)^2$,
\begin{equation}\label{eq:radial-power}
 \boxed{\displaystyle
 \mathscr P(k,t)=\mathscr P(k,0)
 e^{2[\mu-b(k^2-k_\star^2)^2]t}}
 \quad\text{for a.e. }k>0.
\end{equation}

\begin{corollary}[Radial concentration and maximizing wavenumber]
\label{cor:radial}
Assume (H). Then for every $\delta>0$,
\begin{equation}\label{eq:radial-mass}
 \frac{\int_{\{k\ge0:\abs{k-k_\star}\ge\delta\}}
 \mathscr P(k,t)\dd k}
 {\int_0^\infty\mathscr P(k,t)\dd k}
 \longrightarrow0.
\end{equation}
If, in addition, $\mathscr P(\cdot,0)$ has a nonnegative,
nonzero representative in $C_0([0,\infty))$, then every choice
\[
 k_t\in\argmax_{k\ge0}\mathscr P(k,t)
\]
satisfies $k_t\to k_\star$.
If $\mathscr P(k_\star,0)>0$, the estimate
\eqref{eq:rate} holds with $f(\eta_\star)$ and $M$ replaced by
$\mathscr P(k_\star,0)$ and
$\sup_{k\ge0}\mathscr P(k,0)$, respectively.
\end{corollary}
\begin{proof}
Equation \eqref{eq:radial-mass} is precisely
Theorem~\ref{thm:concentration} in polar coordinates.
For the argmax, apply the comparison argument in the proof of
Theorem~\ref{thm:argmax} to the scalar function
$\mathscr P(k,t)$, using \eqref{eq:radial-power}.
The quantitative bound follows by comparing the maximum to the
positive value at $k_\star$, exactly as in
Corollary~\ref{cor:rate}.
\end{proof}

\section{Counterexample and periodic-domain qualification}
\begin{proposition}[Failure without spectral access]\label{prop:counter}
There exist nonzero $A_0\in\mathcal S(\R^d)$ for which
$\argmax P(\cdot,t)$ never approaches $\Sstar$.
\end{proposition}
\begin{proof}
Choose $\delta>0$ and nonzero
$\widehat A_0\in C_c^\infty(\R^d)$ supported in
$\{\xi:\abs{\xi}>k_\star+\delta\}$.
Its inverse Fourier transform is Schwartz. Since the exact
multiplier formula \eqref{eq:ahat} preserves Fourier support,
$P(\xi,t)$ vanishes whenever
$\abs{\xi}\le k_\star+\delta$ for every $t\ge0$.
Each $P(\cdot,t)$ has a positive global maximum and all its
maximizers satisfy $\abs{\xi_t}>k_\star+\delta$.
Thus $\abs{\xi_t}$ cannot converge to $k_\star$.
\end{proof}

\begin{remark}[Flat tori]
On the torus $\T_L^d$ with Fourier lattice
$\xi_n=(2\pi/L)n$, $n\in\Z^d$, the exact coefficientwise law is
\[
 P_n(t)=P_n(0)e^{2\lambda(\xi_n)t}.
\]
The dominant asymptotic modes belong to the set of initially
populated lattice frequencies maximizing $\lambda(\xi_n)$.
Because the lattice is discrete, this need not contain a point with
$\abs{\xi_n}=k_\star$. Therefore convergence of the selected
lattice radius to the continuum $k_\star$ is not automatic for a
fixed torus.
\end{remark}

\section{Conclusion and scope}
The precise E16 result is the exact modal-power identity
\[
 \boxed{P(\xi,t)=P(\xi,0)e^{2\lambda(\xi)t}}
\]
and, under essential spectral access to the maximizing shell,
normalized concentration near $\abs\xi=k_\star$.
Convergence of the pointwise spectral argmax requires an additional
condition ensuring that a maximum exists and controlling the
initial spectral amplitude; $C_0$ regularity is sufficient.
For a scalar radial spectrum, the corresponding statement is
$\argmax_{k\ge0}\mathscr P(k,t)\to k_\star$ under the hypotheses
of Corollary~\ref{cor:radial}.

The theorem does not imply a nonzero whole-space bound state,
nonlinear pattern persistence, phase or directional selection,
or spatial confinement. No such conclusion follows solely from the
linear Fourier multiplier.

\paragraph{Canonical model source.}
The equations correspond to E12--E16 in
\href{https://github.com/rickyjreyes/geometry_of_resonance/blob/main/WCT_FULL_EQUATION_LIST_CORRECTED.md}{WCT Full Equation List (Corrected)}.

\end{document}